\documentclass[12pt,a4paper]{article}
\usepackage[utf8]{inputenc}
\usepackage[T1]{fontenc}
\usepackage[brazil]{babel}
\usepackage{amsmath,amssymb}
\usepackage{booktabs}
\usepackage{geometry}
\geometry{margin=2.5cm}

\title{Primeiro Trabalho de Reconhecimento de Padr\~oes}
\author{}
\date{}

\begin{document}

\maketitle

\section*{Quest\~ao 1}
Obter a fun\c{c}\~ao discriminante baseada na regra de Bayes para um problema de classifica\c{c}\~ao com duas classes equiprov\'aveis:
\[
C_1:\; \mathcal{N}(\mu_1,\Sigma),\qquad
\mu_1=\begin{bmatrix}0\\0\end{bmatrix},\qquad
\Sigma=I
\]
\[
C_2:\; \mathcal{N}(\mu_2,\Sigma),\qquad
\mu_2=\begin{bmatrix}3\\3\end{bmatrix},\qquad
\Sigma=I
\]
com \(P(C_1)=P(C_2)=\frac12\).

\section*{Solu\c{c}\~ao}
A regra de Bayes escolhe a classe que maximiza a probabilidade a posteriori. Como o logaritmo natural \'e uma fun\c{c}\~ao crescente, podemos usar a fun\c{c}\~ao discriminante
\[
g_i(\mathbf{x})=\ln p(\mathbf{x}\mid C_i)+\ln P(C_i).
\]
Para uma Gaussiana multivariada de dimens\~ao \(d\), a densidade condicional \'e
\[
p(\mathbf{x}\mid C_i)=
\frac{1}{(2\pi)^{d/2}|\Sigma_i|^{1/2}}
\exp\left[
-\frac12(\mathbf{x}-\mu_i)^\top\Sigma_i^{-1}(\mathbf{x}-\mu_i)
\right].
\]
Aplicando o logaritmo:
\[
g_i(\mathbf{x})=
-\frac12(\mathbf{x}-\mu_i)^\top\Sigma_i^{-1}(\mathbf{x}-\mu_i)
-\frac{d}{2}\ln(2\pi)
-\frac12\ln|\Sigma_i|
+\ln P(C_i).
\]

Neste problema, \(P(C_1)=P(C_2)\), \(\Sigma_1=\Sigma_2=I\), \(\Sigma^{-1}=I\) e \(|\Sigma|=1\). Portanto, os termos \(\ln P(C_i)\), \(-\frac{d}{2}\ln(2\pi)\) e \(-\frac12\ln|\Sigma_i|\) s\~ao iguais para as duas classes e n\~ao alteram a compara\c{c}\~ao entre elas. Assim, basta comparar
\[
g_i(\mathbf{x})=
-\frac12(\mathbf{x}-\mu_i)^\top(\mathbf{x}-\mu_i).
\]
Expandindo:
\[
g_i(\mathbf{x})=
-\frac12\mathbf{x}^\top\mathbf{x}
+\mu_i^\top\mathbf{x}
-\frac12\mu_i^\top\mu_i.
\]
O termo \(-\frac12\mathbf{x}^\top\mathbf{x}\) tamb\'em \'e comum \`as duas classes. Logo, uma forma equivalente da fun\c{c}\~ao discriminante para cada classe \'e
\[
g_i(\mathbf{x})=
\mu_i^\top\mathbf{x}
-\frac12\mu_i^\top\mu_i.
\]

Substituindo as m\'edias:
\[
\mu_1=\begin{bmatrix}0\\0\end{bmatrix},\quad
\mu_2=\begin{bmatrix}3\\3\end{bmatrix},
\]
temos
\[
g_1(\mathbf{x})=0
\]
e
\[
g_2(\mathbf{x})
=
\begin{bmatrix}3&3\end{bmatrix}
\begin{bmatrix}x_1\\x_2\end{bmatrix}
-\frac12
\begin{bmatrix}3&3\end{bmatrix}
\begin{bmatrix}3\\3\end{bmatrix}
=3x_1+3x_2-9.
\]

A fronteira de decis\~ao ocorre quando \(g_1(\mathbf{x})=g_2(\mathbf{x})\):
\[
3x_1+3x_2-9=0
\;\Longrightarrow\;
x_1+x_2=3.
\]

Assim, uma fun\c{c}\~ao discriminante bin\'aria equivalente \'e
\[
h(\mathbf{x})=x_1+x_2-3.
\]
Regra de decis\~ao:
\[
\begin{cases}
h(\mathbf{x})<0 \Rightarrow \mathbf{x} \in C_1,\\
h(\mathbf{x})>0 \Rightarrow \mathbf{x} \in C_2.
\end{cases}
\]

\paragraph{Interpreta\c{c}\~ao geom\'etrica.}
A fronteira \(x_1+x_2=3\) \'e uma reta com normal \((1,1)\), ortogonal ao vetor \(\mu_2-\mu_1=(3,3)\), e que passa pelo ponto m\'edio entre as m\'edias:
\[
\frac{\mu_1+\mu_2}{2}=\begin{bmatrix}1.5\\1.5\end{bmatrix}.
\]
Isso confirma o resultado esperado para duas Gaussianas com covari\^ancias iguais.

\section*{Quest\~oes 2 a 7 --- Configura\c{c}\~ao original}
Nas quest\~oes seguintes, repetimos o experimento num\'erico descrito no enunciado do trabalho, com duas classes Gaussianas equiprov\'aveis e mesma matriz de covari\^ancia $\Sigma$:
\[
C_1:\;\mathcal{N}(\mu_1,\Sigma),\quad
\mu_1=\begin{bmatrix}0\\0\end{bmatrix},\quad
C_2:\;\mathcal{N}(\mu_2,\Sigma),\quad
\mu_2=\begin{bmatrix}3\\3\end{bmatrix},\quad
\Sigma=I.
\]
Foram geradas $n=1000$ amostras por classe para treinamento (quest\~ao~2) e outras $n=1000$ por classe para teste (quest\~ao~4), com sementes fixadas de forma reproduz\'ivel. A semente global do c\'odigo foi $\texttt{MASTER\_SEED}=2026$, o que produziu semente de treinamento $426$ e semente de teste $89$.

\subsection*{Quest\~ao 2}
Os pontos das duas classes foram plotados com o mesmo estilo do notebook de aula: classe~1 em \texttt{r+}, classe~2 em \texttt{b.}, e a superf\'icie de separa\c{c}\~ao em pontos verdes (\texttt{g.}), sem legenda, eixos \emph{Par\^ametro~1} e \emph{Par\^ametro~2}, grade e t\'itulo. A superf\'icie \'otima de Bayes \'e linear porque $\Sigma_1=\Sigma_2$. A regra linear obtida numericamente coincide com a teoria:
\[
h_{\mathrm{Bayes}}(\mathbf{x})=3x_1+3x_2-9.
\]

\subsection*{Quest\~oes 3 e 4 --- Erro de Bayes}
O erro te\'orico de Bayes para duas Gaussianas com mesma covari\^ancia e priors iguais foi calculado em forma fechada e conferido pela cauda univariada ao longo da normal \`a fronteira; ambos coincidiram em $0{,}0169$. Os erros emp\'iricos (propor\c{c}\~ao de erros de classifica\c{c}\~ao) foram $0{,}0155$ no conjunto de treinamento da quest\~ao~2 e $0{,}0120$ no conjunto de teste da quest\~ao~4.

\subsection*{Quest\~ao 5 --- Classificador Euclidiano}
Usando as m\'edias amostrais do treinamento, a fronteira continua linear, com
\[
h_{\mathrm{Euclidiano}}(\mathbf{x})\approx 2{,}9803\,x_1+2{,}9997\,x_2-9{,}0472.
\]
Como $\Sigma=I$, espera-se forte concord\^ancia com a regra de Bayes; de fato, o erro te\'orico avaliado para essa reta estimada foi $0{,}0170$ e o erro no teste $0{,}0125$.

\subsection*{Quest\~ao 6 --- Bayes ing\^enuo n\~ao param\'etrico}
As densidades marginais em cada dimens\~ao foram estimadas por KDE Gaussiano (fun\c{c}\~ao \texttt{gaussian\_kde} do SciPy); a decis\~ao usa o log-posterior produto das marginais. O erro no teste foi $0{,}0155$. Nesta configura\c{c}\~ao, as vari\'aveis s\~ao independentes sob o modelo gerador ($\Sigma$ diagonal), logo o ing\^enuo aproxima bem o \'otimo.

\subsection*{Quest\~ao 7 --- $k$-NN}
Foram treinados $k$-NN com $k\in\{1,5,11\}$. Os erros no conjunto de teste foram $0{,}0245$, $0{,}0170$ e $0{,}0145$, respectivamente. As matrizes de confus\~ao (linhas = r\'otulo verdadeiro $C_1,C_2$) foram, para $k=1$: $\begin{bmatrix}980&20\\ 29&971\end{bmatrix}$; para $k=5$: $\begin{bmatrix}985&15\\ 19&981\end{bmatrix}$; para $k=11$: $\begin{bmatrix}987&13\\ 16&984\end{bmatrix}$. O valor $k=11$ apresentou o menor erro neste sorteio.

\begin{table}[h]
\centering
\caption{Resumo dos erros no teste --- configura\c{c}\~ao original ($\mu_1=(0,0)^\top$, $\mu_2=(3,3)^\top$, $\Sigma=I$).}
\begin{tabular}{lcc}
\toprule
Classificador & Erro te\'orico & Erro teste \\
\midrule
Bayes \'otimo & $0{,}0169$ & $0{,}0120$ \\
Euclidiano (m\'edias estimadas) & $0{,}0170$ & $0{,}0125$ \\
Bayes ing\^enuo KDE & --- & $0{,}0155$ \\
$k$-NN, $k=1$ & --- & $0{,}0245$ \\
$k$-NN, $k=5$ & --- & $0{,}0170$ \\
$k$-NN, $k=11$ & --- & $0{,}0145$ \\
\bottomrule
\end{tabular}
\end{table}

\section*{Quest\~ao 8 --- Nova configura\c{c}\~ao}
Repetiu-se o mesmo procedimento para
\[
\mu_1=\begin{bmatrix}0\\0\end{bmatrix},\quad
\mu_2=\begin{bmatrix}2\\0\end{bmatrix},\quad
\Sigma=\begin{bmatrix}1&0{,}6\\0{,}6&1\end{bmatrix},
\]
mantendo as mesmas sementes ($426$ e $89$) e $n=1000$ por classe. A fronteira de Bayes permanece linear, com coeficientes num\'ericos
\[
h_{\mathrm{Bayes}}(\mathbf{x})=3{,}1250\,x_1-1{,}8750\,x_2-3{,}1250,
\]
consistentes com $\Sigma^{-1}(\mu_2-\mu_1)$ ap\'os normaliza\c{c}\~ao da reta de decis\~ao. O erro te\'orico de Bayes subiu para $0{,}1056$; no teste, $0{,}0990$. O classificador Euclidiano com m\'edias amostrais deixa de alinhar-se ao \'otimo de Mahalanobis: erro te\'orico da reta estimada $0{,}1599$ e erro no teste $0{,}1645$. O Bayes ing\^enuo por KDE marginal teve erro $0{,}1620$, penalizado pela suposi\c{c}\~ao de independ\^encia quando $\Sigma$ n\~ao \'e diagonal. Para $k$-NN, os erros no teste foram $0{,}1430$ ($k=1$), $0{,}1060$ ($k=5$) e $0{,}0960$ ($k=11$); novamente $k$ maior suaviza a fronteira e melhora a generaliza\c{c}\~ao neste cen\'ario.

\begin{table}[h]
\centering
\caption{Resumo dos erros no teste --- quest\~ao~8 ($\mu_1=(0,0)^\top$, $\mu_2=(2,0)^\top$, $\Sigma$ com correla\c{c}\~ao $0{,}6$).}
\begin{tabular}{lcc}
\toprule
Classificador & Erro te\'orico & Erro teste \\
\midrule
Bayes \'otimo & $0{,}1056$ & $0{,}0990$ \\
Euclidiano (m\'edias estimadas) & $0{,}1599$ & $0{,}1645$ \\
Bayes ing\^enuo KDE & --- & $0{,}1620$ \\
$k$-NN, $k=1$ & --- & $0{,}1430$ \\
$k$-NN, $k=5$ & --- & $0{,}1060$ \\
$k$-NN, $k=11$ & --- & $0{,}0960$ \\
\bottomrule
\end{tabular}
\end{table}

\paragraph{S\'intese.}
Na configura\c{c}\~ao original, classes bem separadas e $\Sigma=I$ fazem Bayes, Euclidiano e ing\^enuo performarem de modo semelhante, com erros baixos. Na quest\~ao~8, maior sobreposi\c{c}\~ao geom\'etrica e correla\c{c}\~ao aumentam o erro de Bayes; o Euclidiano ing\^enuo em rela\c{c}\~ao \`as m\'edias amostrais piora por n\~ao usar $\Sigma^{-1}$, e o Bayes ing\^enuo por marginais perde por modelar incorretamente a depend\^encia. O $k$-NN com $k$ moderado a grande aproxima fronteiras flex\'iveis e compete com o \'otimo quando $k=11$ neste experimento.

\end{document}